\scanpage{104}
\mspage{\deleted{\unsure{48}}49} (\uline{NB} il n'y a pas de page 48)
% archive sheet number: 103

Je résume maintenant ce que je crois avoir compris sur les catégories localisées correspondantes. Il y en a, à équivalence près, \uline{trois} différentes seulement en tout et pour tout, quand on prend les localiseurs ``admis'' dans les catégories, comme je viens de le dire (c'était ça en somme l'intérêt objectif des prescriptions !) : \deleted{\unsure{Précision} $\mathrm{Cat}_{/S}$} \unsure{Précisons mieux|Précision première} \deleted{\unsure{donc}}

1$^\circ$) La $W$-\uline{équivalence} \uline{grossière} donne ``la même'' catégorie localisée pour les 13 catégories du diagramme, i.e. toutes les flèches de transition du diagramme induisent des équivalences pour cette $W$-équivalence \uline{grossière}. Le localisé \deleted{\unsure{correspondant}} sera noté $\boxed{\mathrm{HOT}^{\mathrm{lc}}_W(S)}$, où l'exposant lc signifie : ($W$-types d'homotopie relatifs sur $S$) ``\uline{localement constants}''.

(Il me reste quand même un petit doute pour $\underline{\mathrm{Fibsc}}\,S$ et $\underline{\mathrm{Cofibsc}}\,S$*, si l'équivalence grossière y est bel et bien \unsure{raisonnable|vraisemblable}. On verra en faisant la démonstration.)

% slanted
\margin{Ce \unsure{doute} n'est pas fondé -- \unsure{tout} \deleted{\unsure{va}} \unsure{s'élucide|marche} parfaitement !}

2$^\circ$) La $W$-\uline{équivalence fine à droite} donne le $=$ résultat pour toutes les catégories du diagramme où elle est considérée comme admise, \deleted{\unsure{i.e. dans lesquelles} dans la partie gauche et centrale : 8 catégories $\underbrace{\mathrm{Cat}_{/S},\underline{\mathrm{Fibhot}}\,S,\underline{\mathrm{Parf}}\,S}_{\text{cat. ``centrales''}},\underline{\mathrm{Liss}}\,S,\underline{\mathrm{Fib}}\,S$, \unsure{où coïncide}} \uline{et où elle ne coïncide pas} \uline{avec} \uline{la $W$-équiv. grossière} déjà considérée, donc sur les quatre catégories bordantes gauches

\[
 \mathrm{Cat}_{/S},\ \underline{\mathrm{Liss}}\,S,\ \underline{\mathrm{Fib}}\,S,\ \underline{\mathrm{Fibsc}}\,S.
\]

\margin{\uline{NB} Sur les trois dernières, cette équivalence est la $W$-équiv. par fibres}

Le localisé est noté $\boxed{\mathrm{HOT}_W(S)}$

